\chapter{Pneumothorax}

\label{ch:ch75_pneumothorax}

A pneumothorax occurs when air enters the pleural space, causing partial or complete lung collapse. It may arise spontaneously (especially in tall, thin young adults), after trauma, or as a complication of underlying lung disease. Prompt recognition on imaging is essential.

---

\section{1. Patient Perspective: What a Collapsed Lung Looks Like on Imaging}

If doctors suspect a collapsed lung, you will usually have a \textbf{chest X-ray} first. This is a quick, low-radiation test taken while you stand and take a deep breath.

\begin{itemize}
  \item On a normal chest X-ray, the lungs appear black (filled with air) and the heart and blood vessels appear white.
  \item When a pneumothorax is present, air collects between the lung and the chest wall. The lung edge becomes visible as a thin white line, and the area beyond it is pitch black (no lung markings).
  \item If the pneumothorax is large, the lung may be significantly collapsed.
  \item If you are lying down (e.g., after trauma), the X-ray may be less obvious. Your doctor may then order a \textbf{CT scan} for a more detailed view.
  \item A small pneumothorax may be observed with repeat X-rays, while a large or tension pneumothorax may require immediate treatment with a chest tube.
\end{itemize}

---

\section{2. Imaging Modalities}

\begin{table}[h]
\centering
\begin{tabular}{l l l}
\hline
\textbf{Modality} & \textbf{Advantages} & \textbf{Limitations} \\
\hline
CT Chest & Gold standard; detects occult PTX & Higher dose; not always immediately available \\
Lung ultrasound & Rapid, bedside, high sensitivity & Operator-dependent \\
Supine CXR & Portable for trauma/ICU & Reduced sensitivity; deep sulcus sign \\
Upright PA CXR & Quick, low dose, readily available & May miss small pneumothoraces ($<$ 10\%) \\
\hline
\end{tabular}
\caption{Imaging Modalities for Pneumothorax Detection}
\end{table}

---

\section{3. Imaging Findings}

\subsection{Upright Chest X-Ray}
\begin{itemize}
  \item \textbf{Visceral pleural line}: A thin, sharp white line representing the edge of the visceral pleura, separated from the chest wall by air.
  \item \textbf{Absence of lung markings} peripheral to the visceral pleural line.
  \item \textbf{Deep sulcus sign}: In supine patients, air collects anteriorly and inferiorly, causing an abnormally deep, lucent costophrenic sulcus.
  \item \textbf{Mediastinal shift}: In tension pneumothorax, the mediastinum and trachea shift away from the affected side.
\end{itemize}

\subsection{CT Findings}
\begin{itemize}
  \item CT is the most sensitive modality and can detect pneumothoraces not visible on CXR.
  \item Particularly valuable in \textbf{trauma} (occult pneumothorax), \textbf{bullous disease} (differentiating bullae from pneumothorax), and \textbf{post-procedural} assessment.
  \item CT may reveal associated injuries: rib fractures, haemothorax, pulmonary contusion, or diaphragmatic injury.
\end{itemize}

\subsection{Tension Pneumothorax}
\begin{itemize}
  \item A clinical diagnosis \textbf{not} dependent on imaging.
  \item Imaging findings (if available): complete lung collapse, mediastinal shift to the contralateral side, flattening or inversion of the ipsilateral hemidiaphragm, and compression of the contralateral lung.
  \item \textbf{Do not delay decompression} for imaging in a haemodynamically unstable patient.
\end{itemize}

\begin{pearlbox}[The Deep Sulcus Sign]
On a supine chest X-ray, look for an abnormally deep, lucent costophrenic angle. The deep sulcus sign is the supine equivalent of a visible pneumothorax on an upright film and is easily missed if not specifically sought.
\end{pearlbox}

---

\section{4. Size Classification and Management}

\begin{table}[h]
\centering
\begin{tabular}{l l l}
\hline
\textbf{Size (apex to cupola)} & \textbf{Classification} & \textbf{Typical Management} \\
\hline
2--4 cm & Moderate & Needle aspiration or chest tube \\
$<$ 2 cm & Small & Observation, repeat CXR, consider aspiration \\
$>$ 4 cm or tension & Large/Tension & Immediate decompression + chest tube \\
\hline
\end{tabular}
\caption{Pneumothorax Size and Management}
\end{table}

\section{5. Recurrence Risk and Underlying Disease}

\begin{itemize}
  \item \textbf{Primary spontaneous pneumothorax}: Occurs in otherwise healthy individuals. Recurrence rate $\sim$50\% after first episode. Risk factors include smoking, tall stature, and family history.
  \item \textbf{Secondary spontaneous pneumothorax}: Occurs in patients with underlying lung disease (COPD, cystic fibrosis, asthma, tuberculosis, Marfan syndrome). Higher morbidity and recurrence rate.
  \item \textbf{Traumatic pneumothorax}: May be associated with rib fractures, penetrating injury, or barotrauma (e.g., positive pressure ventilation).
\end{itemize}

\section{6. Bullous Disease Mimics}

Bullae can mimic pneumothorax on CXR. CT differentiates them:
\begin{itemize}
  \item Bullae are rounded, thin-walled air-filled spaces within the lung parenchyma, typically with a visible wall.
  \item Pneumothorax shows air in the pleural space with a visible visceral pleural line and no wall.
  \item CT is essential when the diagnosis is uncertain, especially before chest tube insertion (accidental insertion into a bulla is a complication).
\end{itemize}

---

\section{7. Structured Reporting Example}

\begin{testbox}[Indications:]
22-year-old male, tall build, presenting with sudden-onset right-sided pleuritic chest pain and dyspnoea. Non-smoker. Upright chest X-ray requested.
\end{testbox}
\begin{refbox}[Findings:]
\textbf{Right Hemithorax}: A visceral pleural line is visible in the right upper zone, with no lung markings peripheral to it. The pneumothorax measures approximately 2.5 cm at the apex (moderate size). No mediastinal shift. No tension features.\\
\textbf{Left Hemithorax}: Clear. No pneumothorax, effusion, or consolidation.\\
\textbf{Heart \& Mediastinum}: Normal cardiac silhouette. Trachea is central.\\
\textbf{Bones}: No rib fractures.
\end{refbox}
\begin{pearlbox}[Impression:]
1. Moderate-sized right-sided spontaneous pneumothorax (approximately 2.5 cm).
2. No tension features.
3. Clinical assessment and management per local protocol (observation vs aspiration vs chest tube).
\end{pearlbox}

\clearpage
